\documentclass[10pt,a4paper]{article}
\usepackage[top=1.5cm, bottom=1.5cm, left=1.6cm, right=1.6cm, headheight=14pt, footskip=18pt]{geometry}
\usepackage[utf8]{inputenc}
\usepackage{amsmath,amsfonts,amssymb}
\usepackage{graphicx}
\usepackage{tikz}
\usepackage{circuitikz}
\usetikzlibrary{calc,patterns}
\usepackage[most]{tcolorbox}
\usepackage{enumitem}
\usepackage{fancyhdr}
\usepackage{booktabs}
\usepackage{tabularx}
\usepackage{colortbl}
\usepackage{hyperref}

% --- Palet Warna Resmi UNIB ---
\definecolor{unibblue}{RGB}{0, 32, 96}
\definecolor{unibgold}{RGB}{197, 160, 89}
\definecolor{unibgreen}{RGB}{34, 112, 60}
\definecolor{boxbg}{RGB}{248, 250, 253}
\definecolor{linegray}{RGB}{210, 215, 225}

% --- Pengaturan Header & Footer ---
\pagestyle{fancy}
\fancyhf{}
\lhead{\footnotesize\textbf{\color{unibblue}Universitas Bengkulu} \textbar\ Program Studi S1 Teknik Elektro}
\rhead{\footnotesize\color{gray}Persamaan Diferensial (Genap 2026)}
\lfoot{\footnotesize\color{gray}LKM OBE \textbar\ \textbf{Video}: \url{https://youtu.be/VaRiwphPnaw} \textbar\ \href{https://www.youtube.com/playlist?list=PLS5oOZWXZeTw}{\color{unibblue}\textbf{Playlist}}}
\cfoot{\footnotesize\color{gray}Hal. \thepage\ dari 2}
\rfoot{\footnotesize\color{unibblue}\textbf{Minggu 9: Deret Fourier Trigonometrik \& An...}}
\renewcommand{\headrulewidth}{0.6pt}
\renewcommand{\footrulewidth}{0.4pt}

% --- Custom Box untuk Lembar Jawab ---
\newtcolorbox{answerbox}[1]{%
  colback=white,
  colframe=unibblue!70!black,
  coltitle=white,
  fonttitle=\bfseries\scriptsize,
  colbacktitle=unibblue,
  title={#1},
  arc=2pt,
  left=5pt,right=5pt,top=3pt,bottom=3pt,
  boxrule=0.6pt
}

\begin{document}

% ==================== HALAMAN 1 ====================
\noindent\begin{minipage}[c]{0.68\textwidth}%
  {\large\textbf{\color{unibblue}LEMBAR KERJA MAHASISWA (LKM)}}\\[1mm]
  {\normalsize\textbf{Mata Kuliah: Persamaan Diferensial (2 SKS)}}\\[0.5mm]
  {\small\textbf{Topik Minggu 9:} Deret Fourier Trigonometrik \& Analisis Harmonisa Beban}%
\end{minipage}%
\hfill%
\begin{minipage}[c]{0.30\textwidth}%
  \begin{tcolorbox}[colback=unibblue!5,colframe=unibblue,boxrule=0.6pt,arc=2pt,left=4pt,right=4pt,top=2pt,bottom=2pt]
    \scriptsize
    \textbf{Nama:} \dotfill\\[1.2mm]
    \textbf{NPM:} \dotfill\\[1.2mm]
    \textbf{Kelas/Klp:} \dotfill\\[1.2mm]
    \textbf{Nilai Final:} \hfill \textbf{/ 100}
  \end{tcolorbox}%
\end{minipage}\par

\vspace{1.5mm}

% Kotak Sub-CPMK & Pedoman
\begin{tcolorbox}[colback=boxbg,colframe=unibgold!90!black,boxrule=0.6pt,arc=2pt,left=5pt,right=5pt,top=3pt,bottom=3pt,title=\textbf{\scriptsize Capaian Pembelajaran (Sub-CPMK 5) \& Pedoman Evaluasi OBE}]
  \scriptsize
  \textbf{Sub-CPMK 5:} Mahasiswa mampu merepresentasikan gelombang periodik tak-sinusoidal ke dalam Deret Fourier trigonometrik dan menganalisis distorsi harmonisa (THD) pada sistem kelistrikan sesuai IEEE Std 519.\\
  \textbf{Video Panduan (YouTube):} \url{https://youtu.be/VaRiwphPnaw} \hfill \textbf{Playlist:} \href{https://www.youtube.com/playlist?list=PLS5oOZWXZeTw}{\color{unibblue}\textbf{bit.ly/pd-unib-playlist}}\\\textbf{Alokasi Waktu:} 50--60 Menit \quad\textbar\quad \textbf{Model:} Kolaboratif / Mandiri Terstruktur \quad\textbar\quad \textbf{Bahasa Komputasi:} Julia 1.12+
\end{tcolorbox}

\vspace{1mm}

% Soal C1
\noindent\textbf{1. (C1 - Mengingat \textbar\ Bobot: 10 Poin)}\par\vspace{0.5mm}
{\footnotesize Tuliskan bentuk umum Deret Fourier trigonometrik untuk fungsi periodik $f(t)$ dengan periode $T = 2\pi/\omega_0$, dan tuliskan rumus integral penentuan koefisien $a_0$, $a_n$, dan $b_n$.}
\begin{answerbox}{Ruang Pengerjaan C1 (Definisi Deret Fourier \& Integral Euler-Fourier)}
  \vspace{2.0cm}
\end{answerbox}

\vspace{1mm}

% Soal C2
\noindent\textbf{2. (C2 - Memahami \textbar\ Bobot: 15 Poin)}\par\vspace{0.5mm}
{\footnotesize Jelaskan sifat kesimetrian fungsi genap ($f(-t) = f(t)$), fungsi ganjil ($f(-t) = -f(t)$), dan kesimetrian setengah gelombang (*half-wave symmetry*). Bagaimana kesimetrian ini menyederhanakan perhitungan koefisien Fourier?}
\begin{answerbox}{Ruang Pengerjaan C2 (Pemanfaatan Sifat Kesimetrian Sinyal Listrik)}
  \vspace{2.1cm}
\end{answerbox}

\vspace{1mm}

% Soal C3
\noindent\textbf{3. (C3 - Menerapkan \textbar\ Bobot: 20 Poin)}\par\vspace{0.5mm}
\noindent\begin{minipage}[t]{0.60\textwidth}%
  {\footnotesize Tegangan keluaran inverter fotovoltaik satu fasa berbentuk gelombang persegi simetris ganjil: $v(t) = +V_d$ untuk $0 < t < T/2$ dan $-V_d$ untuk $T/2 < t < T$. Turunkan deret Fourier trigonometrik analitis dari tegangan tersebut.}%
\end{minipage}%
\hfill%
\begin{minipage}[t]{0.37\textwidth}%
  \centering
  \resizebox{0.95\linewidth}{!}{%
  \begin{circuitikz}[scale=0.65, transform shape]
    \draw[->, >=stealth] (-0.2,0) -- (3.5,0) node[right, font=\tiny] {$t$};
    \draw[->, >=stealth] (0,-1.3) -- (0,1.5) node[above, font=\tiny] {$v(t)$};
    \draw[very thick, unibblue] (0,1) -- (1.5,1) -- (1.5,-1) -- (3,-1) -- (3,1) -- (3.3,1);
    \draw[dashed, gray] (1.5,0) node[below, font=\tiny] {$T/2$} -- (1.5,1);
    \draw[dashed, gray] (3,0) node[below, font=\tiny] {$T$} -- (3,-1);
    \node at (-0.4,1) [font=\tiny\bfseries] {$+V_d$};
    \node at (-0.4,-1) [font=\tiny\bfseries] {$-V_d$};
  \end{circuitikz}%
  }%
\end{minipage}\par\vspace{1mm}
\begin{answerbox}{Ruang Pengerjaan C3 (Penurunan Koefisien Fourier Gelombang Inverter)}
  \vspace{2.8cm}
\end{answerbox}

\newpage
% ==================== HALAMAN 2 ====================

% Soal C4
\noindent\textbf{4. (C4 - Menganalisis \textbar\ Bobot: 20 Poin)}\par\vspace{0.5mm}
{\footnotesize Analisis Distorsi Harmonisa Total (\textit{Total Harmonic Distortion} / THD) tegangan yang didefinisikan oleh \textbf{IEEE Std 519}: $\text{THD}_v = \frac{\sqrt{\sum_{n=2}^\infty V_n^2}}{V_1} \times 100\%$. Berdasarkan deret Fourier Soal 3, buktikan bahwa $\text{THD}_v$ gelombang persegi murni bernilai $\approx 48{,}34\%$!}
\begin{answerbox}{Ruang Pengerjaan C4 (Pembuktian Analitis THD Tegangan Standar IEEE 519)}
  \vspace{3.8cm}
\end{answerbox}

\vspace{1.5mm}

% Soal C5
\noindent\textbf{5. (C5 - Mengevaluasi \textbar\ Bobot: 20 Poin)}\par\vspace{0.5mm}
{\footnotesize Evaluasi Fenomena Gibbs: Saat merekonstruksi diskontinuitas loncatan gelombang persegi dengan jumlah harmonisa berhingga ($N \to \infty$), lonjakan (*overshoot*) tidak lenyap melainkan konvergen ke $8{,}95\%$. Jelaskan mengapa penambahan suku harmonisa tinggi tidak dapat menghilangkan lonjakan ini!}
\begin{answerbox}{Ruang Pengerjaan C5 (Evaluasi Fenomena Gibbs \& Konvergensi Titik Diskontinu)}
  \vspace{3.8cm}
\end{answerbox}

\vspace{1.5mm}

% Soal C6
\noindent\textbf{6. (C6 - Merancang \& Komputasi Julia \textbar\ Bobot: 15 Poin)}\par\vspace{0.5mm}
{\footnotesize Rancang skrip program \textbf{Julia} untuk merekonstruksi deret Fourier gelombang persegi hingga harmonisa ke-$N$ ($N=1, 3, 5, 25$) dan memplot perbandingan kurvanya untuk memvisualisasikan fenomena Gibbs.}
\begin{answerbox}{Ruang Pengerjaan C6 (Skrip Rekonstruksi Parsial Deret Fourier Julia)}
  \vspace{3.8cm}
\end{answerbox}

\vspace{1.5mm}

% Tabel Rekapitulasi Skor OBE & Tanda Tangan
\noindent\begin{minipage}[b]{0.65\textwidth}%
  \centering
  \scriptsize
  \begin{tabular}{|c|c|c|c|c|c|c|}
    \hline
    \rowcolor{unibblue}
    \textbf{\color{white}C1} & \textbf{\color{white}C2} & \textbf{\color{white}C3} & \textbf{\color{white}C4} & \textbf{\color{white}C5} & \textbf{\color{white}C6} & \textbf{\color{white}Total} \\
    \rowcolor{unibblue!10}
    10 Pts & 15 Pts & 20 Pts & 20 Pts & 20 Pts & 15 Pts & 100 Pts \\
    \hline
    & & & & & & \\[2.5mm]
    \hline
  \end{tabular}%
\end{minipage}%
\hfill%
\begin{minipage}[b]{0.32\textwidth}%
  \centering
  \scriptsize
  \textbf{Verifikasi Dosen / Asisten:}\\[5mm]
  (\dotfill)
\end{minipage}\par

\end{document}
